\section{Alternative Derivations of \texorpdfstring{$\chi$ and $\xi$}{chi and xi}}\label{app:alt-derivations}

This appendix provides three alternative derivations of the key sufficient
statistics $\chi$ and $\xi$ that yield the same expressions by different
routes. The methods illuminate different aspects of the model structure and
are useful for researchers who wish to extend the framework.

\subsection{Why three methods give the same answer}\label{app:why-three}
The conjecture $\hat c^{H}_{t}=\chi\hat c_{t}+\xi\hat g_{t}$ is the unique
solution to a system of three equilibrium conditions: the \HtM{} budget
constraint, the real-wage aggregation, and the production function
$\hat y_{t}=\hat c_{t}+\hat g_{t}$. Any method that correctly exploits this
system must recover the same $(\chi,\xi)$. The three methods differ in how
they organize the algebra:
(1) \emph{Coefficient matching}: used in Appendix~\ref{app:fullderivation}.
(2) \emph{Matrix system}: expresses the equilibrium conditions as a linear
system and solves by inversion.
(3) \emph{Direct substitution}: substitutes the wage directly into the \HtM{}
budget constraint and solves for $\hat c^{H}_{t}$ as a function of $\hat c_{t}$
and $\hat g_{t}$.

\subsection{Method 1: coefficient matching (main text)}\label{app:method-1}
Reproduced in Appendix~\ref{app:fullderivation} for completeness. The key
algebraic step is matching the coefficients of $\hat c_{t}$ and $\hat g_{t}$
in the equation $(\phi+\sigma\alpha^{H})\hat c^{H}_{t}=KW_{c}\hat c_{t}+[KW_{g}-\sigma\Gamma^{H}-\mu\phi/\lambda]\hat g_{t}$.

\subsection{Method 2: matrix system}\label{app:method-2}
Write the system as $A\mathbf{x}=B\mathbf{g}$ where $\mathbf{x}=(\hat c^{H},\hat c^{S},\hat w,\hat n)'$ and $\mathbf{g}=(\hat c_{t},\hat g_{t})'$:
\[
  A=\begin{pmatrix}
  \phi+\sigma\alpha^{H} & 0 & -K & 0\\
  0 & \phi+\sigma\alpha^{S} & -1 & 0\\
  0 & 0 & 1 & -\phi\\
  -\lambda & -(1-\lambda) & 0 & 1
  \end{pmatrix},
\]
and solve for $\hat c^{H}=f(\hat c,\hat g)$. The $\chi$ and $\xi$ expressions
match those in the main text.

\subsection{Method 3: direct substitution}\label{app:method-3}
The \HtM{} consumption $\hat c^{H}_{t}$ is determined by combining the \HtM{}
budget constraint with the real-wage equation. Direct substitution involves:
(i) solving the wage aggregation for $\hat w_{t}$ as a function of
$\hat c_{t},\hat g_{t},\hat c^{H}_{t}$; (ii) substituting into the \HtM{}
budget \eqref{eq:A-htm-2}; (iii) solving for $\hat c^{H}_{t}$. The result
coincides with Methods 1 and 2.

\subsection{What $\chi$ really is: income cyclicality}\label{app:chi-meaning}
Having derived $\chi$ three ways, it is worth emphasizing the economic
interpretation. $\chi$ is the elasticity of \HtM{} consumption with respect to
aggregate output, holding government spending fixed:
$\chi=\partial\hat c^{H}/\partial\hat y\bigm|_{\hat g=0}$. Since \HtM{} consume
their income, $\chi$ is also the elasticity of \HtM{} income to aggregate
output -- the ``income cyclicality'' terminology. Under countercyclical
inequality ($\chi>1$), \HtM{} income rises more than proportionally with
aggregate output, amplifying the fiscal multiplier through the NK Cross.
